
\documentclass{article}
\usepackage{colortbl}
\usepackage{makecell}
\usepackage{multirow}
\usepackage{supertabular}

\begin{document}

\newcounter{utterance}

\centering \large Interaction Transcript for game `cladder', experiment `full\_v1.5\_default', episode 1084 with llama{-}3.1{-}8b{-}instruct.
\vspace{24pt}

{ \footnotesize  \setcounter{utterance}{1}
\setlength{\tabcolsep}{0pt}
\begin{supertabular}{c@{$\;$}|p{.15\linewidth}@{}p{.15\linewidth}p{.15\linewidth}p{.15\linewidth}p{.15\linewidth}p{.15\linewidth}}
    \# & \multicolumn{2}{c}{Player} && \multicolumn{2}{c}{Game Master} \\
    \hline

    \theutterance \stepcounter{utterance}  
    & & & \multicolumn{4}{p{0.6\linewidth}}{
        \cellcolor[rgb]{0.9,0.9,0.9}{
            \makecell[{{p{\linewidth}}}]{
                \texttt{\tiny{[P1$\langle$GM]}}
                \texttt{You are an expert in causal inference. The following question is not a typical commonsense query, but rather a meticulously designed question created by a professor specializing in causal inference, intended to assess the students' mastery of the course content.} \\
\texttt{Imagine a self{-}contained, hypothetical world with only the following conditions, and without any unmentioned factors or causal relationships: Age has a direct effect on hospital costs and recovery. Hospital costs has a direct effect on recovery. The overall probability of high hospital bill is 51\%. The probability of low hospital bill and recovery is 37\%. The probability of high hospital bill and recovery is 17\%. Is the chance of recovery larger when observing high hospital bill?} \\
\texttt{Start your answer with 'Yes' or 'No', followed by additional reasoning or evidence to support your explanation.} \\
            }
        }
    }
     \\ \\

    \theutterance \stepcounter{utterance}  
    & \multicolumn{4}{p{0.6\linewidth}}{
        \cellcolor[rgb]{0.561,0.737,0.561}{
            \makecell[{{p{\linewidth}}}]{
                \texttt{\tiny{[P1$\rangle$GM]}}
                \texttt{No.} \\
\\ 
\texttt{To answer this question, we need to apply the principles of causal inference, specifically the concept of conditional probability and the causal relationships between variables.} \\
\\ 
\texttt{Let's denote the variables as follows:} \\
\\ 
\texttt{{-} H: High hospital bill (H = 1) or Low hospital bill (H = 0)} \\
\texttt{{-} R: Recovery (R = 1) or No recovery (R = 0)} \\
\\ 
\texttt{We are given the following probabilities:} \\
\\ 
\texttt{{-} P(H = 1) = 0.51 (probability of high hospital bill)} \\
\texttt{{-} P(H = 0, R = 1) = 0.37 (probability of low hospital bill and recovery)} \\
\texttt{{-} P(H = 1, R = 1) = 0.17 (probability of high hospital bill and recovery)} \\
\\ 
\texttt{We want to determine if the chance of recovery is larger when observing a high hospital bill (H = 1). To do this, we need to calculate the conditional probability of recovery given a high hospital bill, P(R = 1 | H = 1).} \\
\\ 
\texttt{Using the definition of conditional probability, we can write:} \\
\\ 
\texttt{P(R = 1 | H = 1) = P(H = 1, R = 1) / P(H = 1)} \\
\\ 
\texttt{Substituting the given values, we get:} \\
\\ 
\texttt{P(R = 1 | H = 1) = 0.17 / 0.51 ≈ 0.333} \\
\\ 
\texttt{Now, let's calculate the probability of recovery given a low hospital bill, P(R = 1 | H = 0). We can use the law of total probability to write:} \\
\\ 
\texttt{P(R = 1 | H = 0) = P(H = 0, R = 1) / P(H = 0)} \\
\\ 
\texttt{Since P(H = 0) = 1 {-} P(H = 1) = 1 {-} 0.51 = 0.49, we get:} \\
\\ 
\texttt{P(R = 1 | H = 0) = 0.37 / 0.49 ≈ 0.755} \\
\\ 
\texttt{Comparing the two conditional probabilities, we see that P(R = 1 | H = 0) > P(R = 1 | H = 1). This means that the chance of recovery is actually larger when observing a low hospital bill, not a high hospital bill.} \\
\\ 
\texttt{Therefore, the answer is No.} \\
            }
        }
    }
    & & \\ \\

    \theutterance \stepcounter{utterance}  
    & & & \multicolumn{2}{p{0.3\linewidth}}{
        \cellcolor[rgb]{0.9,0.9,0.9}{
            \makecell[{{p{\linewidth}}}]{
                \texttt{\tiny{[GM$|$GM]}}
                \texttt{No} \\
            }
        }
    }
    & & \\ \\

    \theutterance \stepcounter{utterance}  
    & & & \multicolumn{2}{p{0.3\linewidth}}{
        \cellcolor[rgb]{0.9,0.9,0.9}{
            \makecell[{{p{\linewidth}}}]{
                \texttt{\tiny{[GM$|$GM]}}
                \texttt{no} \\
            }
        }
    }
    & & \\ \\

    \theutterance \stepcounter{utterance}  
    & & & \multicolumn{2}{p{0.3\linewidth}}{
        \cellcolor[rgb]{0.9,0.9,0.9}{
            \makecell[{{p{\linewidth}}}]{
                \texttt{\tiny{[GM$|$GM]}}
                \texttt{game\_result = WIN} \\
            }
        }
    }
    & & \\ \\

\end{supertabular}
}

\end{document}
